\section{From the CM fibre to the algebraic locus}\label{v7:sec:locus}

\begin{lemma}[Strata of algebraicity]\label{v7:lem:locus}
Let \(f:\mathcal Y\to S\) be smooth and projective over a smooth connected complex algebraic variety, and let \(\gamma\) be a flat rational Hodge class in \(R^{2p}f_*\Q\). The locus \(\Sigma(\gamma)\) where \(\gamma\) is algebraic is a countable union of Zariski closed subsets. It is either all of \(S\) or meagre in the analytic topology.
\end{lemma}
\begin{proof}
For a fixed tuple of Hilbert polynomials, take the product over \(S\) of the corresponding relative Hilbert schemes. It is projective over \(S\). For a fixed tuple of rational coefficients, the condition that the indicated combination of cycle classes equals \(\gamma\) is constant on each connected component: classes of flat families of cycles give flat sections of the pulled back local system. The images of the components satisfying this condition are Zariski closed by properness. Varying the tuples enumerates the rational algebraic representatives. Proper closed subvarieties are nowhere dense in the analytic topology, so the countable union is meagre unless one image is all of \(S\).
\end{proof}

\begin{corollary}[Dense algebraic locus of the mixed link]\label{v7:cor:dense}
In the family of Section~\ref{v7:sec:setup} with its compositum field, \(\Sigma(\xi)\) is dense in the analytic topology. Either \(\Sigma(\xi)\) is the entire curve or it is a countable set.
\end{corollary}
\begin{proof}
The rational \(G\)-action on the period domain carries the special fibre of Theorem~\ref{v7:thm:seed} to isogenous pairs and transports the \(G\)-invariant class \(\xi\). Isogeny pullback and pushforward preserve algebraicity with rational coefficients, so its orbit consists of points of \(\Sigma(\xi)\). The norm-one group has rational points dense in the identity component of its real points. One sees this by the Cayley transform \(x\mapsto (1+x)(1-x)^{-1}\) of the reduced trace-zero subspace, defined over \(F\), which supplies a rational chart about the identity. Rational points of the chart are dense in a real neighbourhood of the identity, and the closed subgroup containing that neighbourhood contains the identity component. The identity component acts transitively on the connected period domain. Passing to the curve yields density. By Lemma~\ref{v7:lem:locus}, a proper \(\Sigma(\xi)\) is a countable union of proper closed subsets of a curve, hence is countable.
\end{proof}

\begin{proposition}[A sufficient bounded-degree condition]\label{v7:prop:bound}
Fix a relatively very ample line bundle on \(\mathcal Y=\mathcal X^2\times_S\mathcal Z\). Suppose that, at infinitely many distinct fibres of the Shimura curve \(S\), \(\xi\) lies in the rational span of codimension-two subvarieties of degree at most one fixed integer \(D_0\). Then \(\xi\) is algebraic at every fibre. No bound on the denominators of the coefficients is needed.
\end{proposition}
\begin{proof}
Let \(b=\dim_\Q H^4(\mathcal Y_s,\Q)\). At each of the chosen fibres, a spanning collection for \(\xi\) can be reduced to at most \(b\) classes. Parameterize ordered tuples of at most \(b\) effective cycles, each of degree in \(\{1,\ldots,D_0\}\), by the corresponding products of relative Chow varieties. These parameter spaces are proper over \(S\) and have finitely many components \cite{Ful98}. On each connected component the span of the classes is flat, so containing the fixed flat class \(\xi\) is a componentwise condition. The images of the admissible components are therefore a finite collection of closed subsets of \(S\). They contain the infinitely many chosen fibres, so one of the images is all of \(S\). This gives the claim. The span condition uses only the classes of the supports and places no restriction on rational coefficients.
\end{proof}

\subsection{The missing deformation}

The dense orbit in Corollary~\ref{v7:cor:dense} gives no uniform bound of the kind in Proposition~\ref{v7:prop:bound}. At full-group fibres, the joint Hodge group has derived group the common image of \(G\), since it projects onto the Hodge group of \(X\) and the two representations use the same left circle. The \(G\)-representations \(H^1(X)\) and \(H^1(Z)\) have no common irreducible summand. The first has the tensor product of all three standard factors and the second has only single standard factors. Therefore
\[
 \Hom(X,Z)\otimes\Q=0
\]
at such fibres. The isogeny retraction used to build the cycle at \(c\) cannot be used as a family of homomorphisms through a full-group fibre.

This also identifies a limit of the divisor construction. When the cross homomorphisms vanish, the divisors of \(X^2\times Z\) are sums of classes from \(X^2\) and from \(Z\). A product of divisors in the mixed K\"unneth component \(H^1(X)\otimes H^1(X)\otimes H^2(Z)\) has its first two factors in the span of the Poincar\'e tensor. That span is the trivial representation and is disjoint from \(T_X\). Thus divisor products cannot supply \(\xi\) on a full-group fibre even though divisors suffice on every power of \(Z\).

\subsection{The first infinitesimal test}\label{v7:sec:firstorder}

\begin{proposition}\label{v7:prop:firstorder}
At the matched CM fibre, no nonzero rational homomorphism \(X_c\to Z_c\), or \(Z_c\to X_c\), extends as a flat Hodge morphism over the first-order deformation in the nonzero tangent direction of the Shimura curve.
\end{proposition}
\begin{proof}
Work over an algebraic splitting field for the torus at \(c\). Its action on the pair is the graph torus of Theorem~\ref{v7:thm:seed}. A rational Hodge morphism \(f:H^1(X_c)\to H^1(Z_c)\) commutes with that torus. The infinitesimal period map for the varying left circle is a nonzero element \(N_3\) of the negative root line in the third \(\mathfrak{sl}_2\) factor. It has Hodge type \((-1,1)\) on both weight-one representations. On \(D_1,D_2\) its action is zero, and on \(D_3\) it is the left lower-root action.

Extending \(f\) as a flat Hodge morphism to the first-order deformation requires \([N_3,f]=0\) on \(H^{1,0}\). This commutator has type \((-1,1)\), so it is already zero on \(H^{0,1}\), whose image would have type \((-1,2)\). Thus \(f\) intertwines the entire root line. Since \(f\) is rational, it intertwines every Galois conjugate of the line. The embeddings of \(F\) are transitive, and conjugation in the common imaginary quadratic field reverses all root signs. These conjugates are the six root lines \(\pm2e_i\); they generate \(\operatorname{Lie}(G)\). It follows that \(f\) intertwines \(G\). But \(V_1\otimes V_2\otimes V_3\) and the sum of the single-standard representations \(D_i\) have no common constituent, so \(f=0\). The argument in the reverse direction is identical.
\end{proof}

The proposition concerns the weight-one morphisms used in the construction. It does not give a first-order obstruction to the final degree-four class \(\xi\), which is flat and of Hodge type along the curve.

\subsection{Representatives and semiregularity}

Semiregularity provides a route to studying deformations of the Chern character of a perfect complex. The theorem identifies the image under semiregularity of an obstruction with the obstruction to preserving the Chern character as a Hodge class \cite{Pri24}. When the character stays Hodge, the obstruction is therefore in the kernel. Its vanishing requires an additional condition on the representative, such as injectivity of the semiregularity map. An algebraic class at the CM fibre gives a rational \(K\)-theory class, but it does not determine these self-extension groups.

\begin{proposition}[A representative with pure Chern character]\label{v7:prop:adams}
There is a perfect complex \(E\) on \(Y_c=X_c^2\times Z_c\) and an integer \(N>0\) with
\[
 \operatorname{ch}_2(E)=N\xi_c,\qquad
 \operatorname{ch}_j(E)=0\quad(j\ne2)
\]
in rational cohomology. All the components of this Chern character extend as flat Hodge classes along the curve.
\end{proposition}
\begin{proof}
Choose a rational codimension-two cycle representing \(\xi_c\) by Theorem~\ref{v7:thm:seed}. Taking the same rational combination of its structure sheaves gives a class \(v\in K_0(Y_c)_\Q\) with \(\operatorname{ch}_2(v)=\xi_c\). The second Adams operation \(\Psi^2\) satisfies \(\operatorname{ch}_j(\Psi^2v)=2^j\operatorname{ch}_j(v)\) \cite{Ful98}. Since \(\dim Y_c=14\), apply the rational polynomial operator
\[
 e_2=\prod_{\substack{0\le j\le14\\j\ne2}}
              \frac{\Psi^2-2^j}{4-2^j}.
\]
On the \(j\)-th component of the Chern character it acts as the indicator of \(j=2\). Clear denominators in \(e_2v\), and write the resulting integral \(K_0\) class as a difference of vector bundles, represented by a perfect complex. Its cohomological character is the displayed one. The nonzero component extends by \(N\xi\) and the remaining components extend by zero.
\end{proof}

This constructs the character without a deformation estimate for its representative. In the first-order direction with Kodaira--Spencer class \(\kappa\), the obstruction is the contraction of the Atiyah class with \(\kappa\), in \(\operatorname{Ext}^2(E,E)\). Its semiregularity image vanishes because the full Chern character stays Hodge. The calculation therefore requires vanishing of the obstruction inside the kernel, or injectivity on the obstruction classes, rather than only preservation of the character.

\begin{lemma}[Direct sums and the trace map]\label{v7:lem:directsum}
Let \(E_1,E_2\) be perfect complexes on a smooth projective variety. The semiregularity map of \(E_1\oplus E_2\) vanishes on the two off-diagonal summands
\(\operatorname{Ext}^2(E_1,E_2)\) and \(\operatorname{Ext}^2(E_2,E_1)\).
For \(E\oplus E\), its kernel contains
\(\operatorname{Ext}^2(E,E)\otimes\mathfrak{sl}_2\).
\end{lemma}
\begin{proof}
The Atiyah class of a direct sum is block diagonal. Each component of the semiregularity map is the trace of the product of a self-extension with a power of that class. Multiplying by an off-diagonal matrix leaves it off diagonal, with trace zero. For two identical summands the formula is the semiregularity map on \(E\) multiplied by the ordinary matrix trace, proving the second assertion.
\end{proof}

In particular \(L\oplus L\) is not semiregular on an abelian variety of dimension at least two, since \(\operatorname{Ext}^2(L,L)=H^2(\mathcal O)\ne0\). Direct sums of known objects therefore do not provide injectivity by formal additivity. A cone changes the extension groups and the differential of the complex. Its Chern character is additive, while the required deformation calculation depends on that differential.

For the mixed link it would suffice to construct a representative at \(c\) whose relative deformation space dominates the curve and whose entire Chern character remains Hodge. A dominating family gives algebraicity on an open set and Lemma~\ref{v7:lem:locus} then gives it everywhere. Neither such a representative nor the bounded-degree input is established here.

These results prove the seed statement and the algebraicity statements for the sixfold. They reduce the Mumford-square problem at a full-group fibre to a specific mixed cycle or one exceptional cycle on the square. Even a proof of that Mumford-square statement would address only one part of \(\Ftwo\). The statement \(\Fthree\) for arbitrary nonabelian varieties requires separate algebraic correspondences. The two universal statements, and hence the full Hodge conjecture, remain open.
